% !TeX root = ../../Book4.tex
% 纲要标题：### 14.3 导出逆极限
% 纲要标签：b4:sec:1303
% 纲要位置：计划纲要.md，第 2811 行。

\section{导出逆极限}
\label{b4:sec:1303}

逆极限的高次导出需要对整张相容图表取内射消解。可数塔可以用乘积差映射压缩这一计算，一般指标范畴则可能保留更多高次信息。

% 纲要内容（仅作写作计划，尚非教材正文）：
%
% * R^i lim 的定义
% * 可数塔的乘积差映射
% * 可数塔与一般指标范畴的高阶差异
%

\subsection{\texorpdfstring{逆极限的导出函子 \(R^i\!\varprojlim\)}{逆极限的导出函子}}
\label{b4:subsec:1303:01}

设 \(R\) 为环，\(\mathcal I\) 为小范畴。模图表的极限是加性函子，也是常值图表函子的右伴随，因而左正合。由\cref{b4:prop:diagram-enough-injectives}，可以在图表范畴中取内射消解
\[
0\longrightarrow M\longrightarrow I^0\longrightarrow I^1\longrightarrow\cdots.
\]
定义\term[daochu-nijixian]{导出逆极限}[derived inverse limit]
\[
R^q\!\varprojlim_{\mathcal I} M
\coloneqq H^q\bigl(\varprojlim_{\mathcal I} I^\bullet\bigr),\qquad q\ge0.
\]
\symidx{\(R^q\varprojlim\)：逆极限的右导出函子}
消解比较与泛性已在\cref{b4:thm:resolution-comparison}及\cref{b4:thm:derived-universal}建立。这里消解的对象是整张图表；在各个位置分别任取内射消解而不构造相容箭头，不能直接代入此定义。

\ProseParagraphBreak
零次为通常极限。一个图表短正合列产生长正合列，其第一个连接态射衡量相容提升失败。以下对可数塔把这套抽象构造化为只有两项的具体复形。

\subsection{可数塔的乘积差映射}
\label{b4:subsec:1303:02}

\begin{theorem}\label{b4:thm:countable-derived-limit}
设 \(R\) 为环，\(M=(M_n,u_n)\) 为左 \(R\)-模可数逆向塔。令
\[
\Delta_M:\prod_{n\ge0}M_n\longrightarrow\prod_{n\ge0}M_n,
\qquad \Delta_M(x)_n=x_n-u_nx_{n+1}.
\]
则自然地有
\[
\varprojlim M_n=\ker\Delta_M,\qquad
R^1\!\varprojlim M_n\simeq\operatorname{coker}\Delta_M,
\qquad R^q\!\varprojlim M_n=0\quad(q\ge2).
\]
若每个 \(u_n\) 满射，则 \(R^1\!\varprojlim M_n=0\)，且极限到每个 \(M_n\) 的投影满射。
\end{theorem}
\begin{proof}
核的等式正是相容性定义。暂记 \(T^0(M)=\ker\Delta_M\)、\(T^1(M)=\operatorname{coker}\Delta_M\)，并令 \(T^q(M)=0\) 对 \(q\ge2\)。塔的态射与各过渡映射交换，故其诱导的积映射与 \(\Delta\) 交换，从而给出 \(T^0,T^1\) 上的映射。

对任意塔短正合列 \(0\to A\to B\to C\to0\)，逐层取直积得到两行短正合列，并以 \(\Delta_A,\Delta_B,\Delta_C\) 作竖直箭头。模的直积正合，蛇引理给出
\[
\begin{gathered}
0\longrightarrow T^0(A)\longrightarrow T^0(B)\longrightarrow T^0(C)
\xrightarrow{\partial}T^1(A)\\
\longrightarrow T^1(B)\longrightarrow T^1(C)\longrightarrow0.
\end{gathered}
\]
蛇引理的自然性保证 \(\partial\) 与短正合列的态射相容；在后面接上为零的 \(T^q\)，这就得到上同调 \(\delta\)-函子 \(T^\bullet\)。

若过渡映射均满，给定 \(y=(y_n)\)，任选 \(x_0\)，逐次选 \(x_{n+1}\) 使 \(u_nx_{n+1}=x_n-y_n\)，便得 \(\Delta(x)=y\)。所以 \(T^1(M)=0\)。为证明极限投影满，固定某个位置的元素，向前用过渡映射确定较小指标的值，向后逐次取原像即可。

对任意塔 \(M\)，\cref{b4:prop:diagram-enough-injectives}提供一个塔范畴中的单态射 \(\iota:M\hookrightarrow I\)，其中 \(I\) 内射且各过渡映射满射。刚才的计算给出 \(T^1(I)=0\)，所以 \(T^1(\iota):T^1(M)\to T^1(I)\) 为零。这对每个 \(M\) 都成立，正是 \(T^1\) 的可消去性。各 \(T^q\) 在 \(q\ge2\) 时已经为零，也自动可消去。

因此，\cref{b4:thm:derived-universal}的消去性判据说明 \(T^\bullet\) 是泛 \(\delta\)-函子。另一方面，内射消解定义的 \(R^\bullet\!\varprojlim\) 也是泛 \(\delta\)-函子，且两者的零次均为 \(\varprojlim\)。零次恒等由泛性唯一延伸为两个方向的 \(\delta\)-函子态射。它们的两个复合在零次都是恒等，因此由延伸的唯一性，在所有次数也都是恒等。这便得到保持连接态射的典范自然同构
\[
T^q(M)\xrightarrow{\ \sim\ }R^q\!\varprojlim M
\qquad(q\ge0),
\]
从而同时证明一次余核计算公式与高次消失。
\end{proof}

以后对可数模塔把 \(R^1\!\varprojlim\) 简写成 \(\varprojlim{}^1\)。对短正合塔列，其长正合列只有
\[
\begin{gathered}
0\to\varprojlim A_n\to\varprojlim B_n\to\varprojlim C_n
\xrightarrow{\partial}\varprojlim{}^1 A_n\\
\longrightarrow\varprojlim{}^1B_n\longrightarrow\varprojlim{}^1C_n\longrightarrow0.
\end{gathered}
\]
\cref{b4:prop:compatible-lift-obstruction}正是这里 \(\partial\) 的元素计算。

\begin{example}\label{b4:exm:derived-limit-adic-quotient}
设 \(p\) 为素数。对\subsecref{b4:subsec:1302:01}的整数塔，常值塔 \(\mathbb Z\) 及约化塔 \(\mathbb Z/p^{n+1}\) 的过渡映射都满。上列因而给出
\[
\varprojlim{}^1_n p^{n+1}\mathbb Z\simeq\mathbb Z_p/\mathbb Z\ne0.
\]
把 \(p^{n+1}\mathbb Z\) 用 \(p^{n+1}a\leftrightarrow a\) 识别为 \(\mathbb Z\)，其过渡映射变成乘 \(p\)。所以这也算出了常取值、非恒等过渡塔 \(\mathbb Z\xleftarrow{p}\mathbb Z\xleftarrow{p}\cdots\) 的一次导出极限。
\end{example}

\subsection{可数塔与一般指标范畴的高阶差异}
\label{b4:subsec:1303:03}

可数塔的箭头由相邻过渡唯一确定，允许用一族差值 \(x_n-u_nx_{n+1}\) 记录全部关系，所以两项乘积复形已经足够。一般图表还可能具有额外的复合关系，产生更高次的导出极限。

\begin{example}\label{b4:exm:general-index-higher-limit}
取仅有一个对象、态射群为 \(C_2\) 的范畴 \(BC_2\)。交换群值函子等价于左 \(\mathbb Z[C_2]\)-模，极限等价于不变量函子。因此其导出极限是\cref{b4:def:group-cohomology}的群上同调。对平凡模 \(M=\mathbb F_2\)，\cref{b4:thm:cyclic-group-cohomology}的周期公式给出
\[
R^q\!\varprojlim_{BC_2}M=H^q(C_2,\mathbb F_2)=\mathbb F_2
\qquad(q\ge0).
\]
这个有限指标范畴已经有任意高次非零导出极限；它不是一座塔，也不是有向偏序集。
\end{example}

对任意集合指标的偏序集，若值域是满足 AB4\(^*\) 的交换范畴，即小积存在且正合，就可以构造计算高次相容性的多项复形；若值域还有足够多内射对象，该复形的上同调便与逆极限的右导出函子典范同构。这一构造及其条件见 Weibel\cite[Vista 3.5.12, p. 86]{Weibel1994HomologicalAlgebra}。指标范畴的复合关系与值域的积正合性，共同决定这种复形能否计算导出极限。
